\def\EDF{\text{EDF}_\alpha}
\def\OPT{\text{OPT}}

\chapter{Analysis of the $\EDF$ algorithm}

A result from Chin \emph{et al.} \cite{CHIN2006255} states that the $\EDF$
algorithm is $\phi$-competitive on $3$-bounded instances. Therefore, it cannot
be worse than $\phi$-competitive on $(2,3)$-bounded instances, which are a
subset of $3$-bounded instances. In this chapter, we show that the competitive
ratio cannot be further improved on $(2,3)$-bounded instances.

\section{Preliminaries}

\begin{algorithm}
\caption{The $\EDF$ Algorithm}
\begin{algorithmic}[1]

\While{there are pending packets in the buffer}
\State $h \gets \text{heaviest pending packet}$
\State $d \gets \text{pending packet with the earliest deadline satisfying } \alpha w_d \geq w_h$
\State $e \gets \text{heaviest pending packet with the same deadline as } d$
\State schedule $e$
\EndWhile
\end{algorithmic}
\end{algorithm}

In this chapter we define $\beta$ as $\sqrt{\alpha}$.

We assume that the additive constant $a$ in the competitive ratio is $0$. This
assumption can be easily removed since the $\EDF$ algorithm does not change
scheduled packets if the weights of all packets are scaled by the same factor.
Thus, we can multiply the weight of each packet by $x \gg a$, making
the additive constant negligible.

\section{Results}

\begin{lemma}\label{lem:edf-I1}
There exists a $(2,3)$-bounded instance where the competitive ratio of $\EDF$
is at least $1 + \frac{1}{\alpha} - \varepsilon$ for all parameters $\alpha \geq 1$
and $\varepsilon > 0$.
\end{lemma}

\begin{proof}
Let $\delta$ be equal to $\varepsilon \alpha$. Since the competitive ratio cannot
be less than $1$, the lemma holds trivially when $\frac{1}{\alpha} \leq
\varepsilon$. Therefore, we can assume that $\frac{1}{\alpha} - \varepsilon \in (0,1)$.

\begin{gather*}
\varepsilon \alpha = \delta > \delta (\frac{1}{\alpha} - \varepsilon) \\
0 > \frac{\delta}{\alpha} - \varepsilon(\alpha + \delta) \\
1 > \frac{\delta + \alpha}{\alpha} - \varepsilon(\alpha + \delta) \\
\frac{1}{\alpha + \delta} > \frac{1}{\alpha} - \varepsilon
\end{gather*}

We will construct an instance $I_1$ (as shown in Figure \ref{img:edf-I1}). In
the first step, the adversary releases two packets of weight $1$ with span $2$
and one packet of weight $\alpha + \delta$ with span $3$. Because $\alpha +
\delta > \alpha \cdot 1$, the $\EDF$ algorithm schedules the heaviest packet. In
the next step, the adversary releases a single packet with weight $\alpha +
\delta$ and span $3$. The $\EDF$ algorithm again schedules the heaviest packet,
resulting in a gain of $W_{\EDF} = 2(\alpha + \delta)$.

The optimal algorithm can schedule all packets, resulting in a gain of $W_{\OPT}
= 2 + 2(\alpha + \delta)$. The competitive ratio is:

\[
\frac{W_{\OPT}}{W_{\EDF}} = \frac{2 + 2(\alpha + \delta)}{2(\alpha + \delta)} = 1 + \frac{1}{\alpha + \delta}
\]

By the choice of $\delta$ we have that $1 + \frac{1}{\alpha + \delta} > 1 +
\frac{1}{\alpha} - \varepsilon$, which concludes this proof.
\end{proof}

\begin{figure}[h]\centering
\includegraphics{img/edf-i1}
\caption{Instance $I_1$ used in Lemma \ref{lem:edf-I1}.}
\label{img:edf-I1}
\end{figure}

\begin{lemma}\label{lem:edf-I2}
There exists a $(2,3)$-bounded instance where the competitive ratio of $\EDF$
is at least $2 - \frac{2 - \beta}{\beta^2 - \beta + 1} - \varepsilon$ for all
parameters $\alpha \geq 1$ and $\varepsilon > 0$.
\end{lemma}

\begin{proof}
To begin the instance $I_2$, the adversary releases two packets of weight $1$ with
span $2$ and one packet of weight $\beta$ with span $3$. At step $t$, where $2
\leq t \leq n$, the adversary releases one packet of weight $\beta^t$ with span
$3$. To end the instance, at step $n+1$, the adversary releases two packets of weight
$\beta^{n+1}$ with span $2$.

In the first step, the $\EDF$ algorithm has pending packets of weights $1$ and
$\beta$, so it schedules the most urgent packet. In step $2 \leq t \leq n$, the
pending packets of $\EDF$ have weights $\beta^{t-2}$, $\beta^{t-1}$, and
$\beta^t$. Since $\alpha \beta^{t-2} \geq \beta^t$, $\EDF$ again schedules the
most urgent packet. In step $n+1$, the pending packets are of weights
$\beta^{n-1}$, $\beta^n$, and $\beta^{n+1}$, so $\EDF$ schedules the most urgent
packet. In the last step $n+2$, the algorithm schedules the heaviest packet of
weight $\beta^{n+1}$. The total gain of the $\EDF$ algorithm is:

\[
W_{\EDF} = 1 + \sum_{i = 0}^{n-1} {\beta^i} + \beta^{n+1} = 1 + \frac{\beta^n - 1}{\beta-1} + \beta^{n+1}
\]

We now compute the limit of $\frac{W_{\EDF}}{\beta^n}$ as $n$ approaches $\infty$:

\begin{equation}\label{I2:eq1}
\lim_{n \to \infty} \frac{W_{\EDF}}{\beta^n} = \lim_{n \to \infty}
\frac{1}{\beta^n} + \frac{1}{\beta-1} - \frac{1}{\beta^n(\beta - 1)} + \beta =
\beta + \frac{1}{\beta - 1}
\end{equation}

The optimal algorithm can take the heaviest packet in every step, resulting in a gain of:

\[
W_{\OPT} = \sum_{i = 1}^{n+1} {\beta^i} + \beta^{n+1} = \frac{\beta^{n+2} - \beta}{\beta - 1} + \beta^{n+1}
\]

We again compute the limit of $\frac{W_{\OPT}}{\beta^n}$ as $n$ approaches $\infty$:

\begin{equation}\label{I2:eq2}
\lim_{n \to \infty} \frac{W_{\OPT}}{\beta^n} = \frac{\beta^2}{\beta - 1} - \frac{1}{\beta^{n-1}(\beta - 1)} + \beta = \frac{\beta^2}{\beta - 1} + \beta
\end{equation}

Combining \ref{I2:eq1} and \ref{I2:eq2}, we compute the limit of the competitive
ratio:

\[
\lim_{n \to \infty} \frac{W_{\OPT}}{W_{\EDF}} = \frac{\frac{\beta^2}{\beta - 1} + \beta}{\beta + \frac{1}{\beta - 1}}
= \frac{2\beta^2 - \beta}{\beta^2 - \beta + 2}
= 2 - \frac{2 - \beta}{\beta^2 - \beta + 1}
\]

Therefore for all $\varepsilon > 0$ there exists $n$ such that
$\frac{W_{\OPT}}{W_{\EDF}} \geq 2 - \frac{2 - \beta}{\beta^2 - \beta + 1} -
\varepsilon$.
\end{proof}

\begin{figure}[h]\centering
\includegraphics{img/edf-i2.pdf}
\caption{Instance $I_2$ used in Lemma \ref{lem:edf-I2}.}
\label{img:edf-I2}
\end{figure}

\begin{lemma}\label{lem:edf-I3}
There exists a $(2,3)$-bounded instance where the competitive ratio of $\EDF$
is at least $1 + \frac{3\beta^2 - 2\beta - 1}{2\beta^4 - 2\beta^3 + 1} -
\varepsilon$ for all parameters $\alpha \geq 1$ and $\varepsilon > 0$.
\end{lemma}

\begin{proof}
The beginning of instance $I_3$ is the same as in Lemma \ref{lem:edf-I2}. To end
the instance, in step $n+1$, the adversary releases two packets of weight
$\beta^{n}$ with span $2$ and one packet of weight $\beta^{n+2} + 1$ with
span $3$. In step $n+2$, the adversary releases a single packet of weight
$\beta^{n+2} + 1$ with span $3$.

In the first $n$ steps, the $\EDF$ algorithm behaves exactly the same as on
instance $I_2$ \ref{lem:edf-I2}, taking the most urgent packets. At step $n+1$,
the $\EDF$ algorithm has pending packets of weights $\beta^{n-1}$, $\beta^n$ and
$\beta^{n+2} + 1$. Since $\alpha \beta^n < \beta^{n+2} + 1$, the
algorithm schedules the heaviest packet. At step $n+2$, the algorithm has
pending packets of weights $\beta^n$ and $\beta^{n+2} + 1$, so it again
schedules the heaviest packet. This results in a gain of:

\[
W_{\EDF} = 1 + \sum_{i = 0}^{n-2} \beta^i + 2(\beta^{n+2} + 1) = 1 + \frac{\beta^{n-1} - 1}{\beta - 1} + 2(\beta^{n+2} + 1)
\]

Taking the limit of $\frac{W_{\EDF}}{\beta^{n-1}}$ as $n$ approaches $\infty$ gives us:

\begin{equation}\label{I3:eq1}
\lim_{n \to \infty} \frac{W_{\EDF}}{\beta^{n-1}} =
\lim_{n \to \infty} \frac{1}{\beta^{n-1}} + \frac{1}{\beta - 1} -
\frac{1}{\beta^{n-1}(\beta - 1)} + 2\beta^3 + \frac{2}{\beta^{n-1}} =
2 \beta^3 + \frac{1}{\beta + 1}
\end{equation}

The optimal algorithm can take the heaviest packet in each step, resulting in a gain of:

\[
W_{\OPT} = \sum_{i=1}^n \beta^i + 2\beta^n + 2(\beta^{n+2} + 1) =
\frac{\beta^{n+1} - \beta}{\beta - 1} + 2\beta^n + 2(\beta^{n+2} + 1)
\]

We again compute the limit of $\frac{W_{\OPT}}{\beta^{n-1}}$ as $n$ approaches $\infty$:

\begin{equation}\label{I3:eq2}
\lim_{n \to \infty} \frac{W_{\OPT}}{\beta^{n-1}} =
\lim_{n \to \infty} \frac{\beta^2}{\beta - 1} -
\frac{1}{\beta^{n-2}(\beta - 1)} + 2\beta + 2\beta^3 + \frac{2}{\beta^{n-1}} =
2\beta^3 + 2\beta + \frac{\beta^2}{\beta-1}
\end{equation}

Combining \ref{I3:eq1} and \ref{I3:eq2}, we compute the limit of the competitive
ratio:

\[
\lim_{n \to \infty} \frac{W_{\OPT}}{W_{\EDF}} =
\frac{2\beta^3 + 2\beta + \frac{\beta^2}{\beta-1}}{2\beta^3 + \frac{1}{\beta - 1}} =
\frac{2\beta^4 - 2\beta^3 + 3\beta^3 - 2\beta}{2\beta^4 - 2\beta^3 + 1} =
1 + \frac{3\beta^2 - 2\beta - 1}{2\beta^4 - 2\beta^3 + 1}
\]

Therefore for all $\varepsilon > 0$ there exists $n$ such that
$\frac{W_{\OPT}}{W_{\EDF}} \geq 1 + \frac{3\beta^2 - 2\beta - 1}{2\beta^4 - 2\beta^3 + 1} - \varepsilon$.
\end{proof}

\begin{figure}[h]\centering
\includegraphics{img/edf-i3.pdf}
\caption{Instance $I_3$ used in Lemma \ref{lem:edf-I3}.}
\label{img:edf-I3}
\end{figure}

\begin{thm}
The competitive ratio of $\EDF$ algorithm on $(2,3)$-bounded instances is at
least $\phi - \varepsilon$ for all $\varepsilon > 0$.
\end{thm}
\begin{proof}
Based on the known parameter $\alpha$, the adversary can choose one of the three
instances from Lemmas \ref{lem:edf-I1}, \ref{lem:edf-I2}, and \ref{lem:edf-I3}
which maximizes the competitive ratio.

\xxx{TODO}
\end{proof}
